\chapter[GLMs and optimization]{Exponential families, GLMs, and constrained optimization}
\label{ch:glm-optimization}
The Poisson model in \path{poisson_spike_glm.m}
can be derived from a probabilistic model, rather than treated
as an unexplained numerical recipe.

\section{Poisson counts as an exponential family}
For count $Y\in\{0,1,\ldots\}$ and positive mean $\mu$,
\[
p(y\mid\mu)=\exp\bigl(y\log\mu-\mu-\log(y!)\bigr).
\]
Its natural parameter is $\eta=\log\mu$. For a row of
covariates $x_i^\top$, exposure duration $\Delta_i>0$,
and regression coefficients $\beta$, define
$\mu_i=\Delta_i\exp(x_i^\top\beta)$.
In matrix form with design $X$,
\[
\ell(\beta)=\sum_i[y_i(x_i^\top\beta+\log\Delta_i)
 -\Delta_i e^{x_i^\top\beta}-\log(y_i!)].
\]
\begin{proposition}[Concavity of the Poisson log likelihood]
For fixed nonnegative counts and strictly positive exposures,
\[
\nabla\ell=X^\top(y-\mu),\qquad
\nabla^2\ell=-X^\top W X,\quad W=\diag(\mu_i)\succ0.
\]
Thus $\ell$ is concave; if $X$ has full column rank, it is
strictly concave at every finite coefficient vector.
\end{proposition}
\begin{proof}
Differentiate the exponential using
\[\frac{\partial\mu_i}{\partial\beta}=\mu_i x_i.\]
For nonzero $v$, $-v^\top\nabla^2\ell v
=\sum_i\mu_i(x_i^\top v)^2$, positive under full rank.
\end{proof}
\emph{Strict concavity is not existence of a finite maximizer.}
For all-zero counts and an unpenalized intercept, likelihood
increases as the intercept tends to $-\infty$.

\section{IRLS as Fisher scoring}
Set $\eta=X\beta+\log\Delta$ and $\mu=e^\eta$.
Newton scoring solves, at the current iterate,
\[
(X^\top W X)\,\delta=X^\top(y-\mu),
\qquad \beta_{\mathrm{new}}=\beta+\delta.
\]
This is equivalent to weighted least squares on the working
response $z=\eta+(y-\mu)/\mu$, with the exposure offset
kept fixed. For the canonical Poisson family, the observed
negative Hessian agrees with expected Fisher information.
A quadratic ridge penalty adds $-\lambda P\beta$ to the score
and $\lambda P$ to the negative Hessian. If the algorithm
\emph{clips} $\eta$, or lacks line search, the numerical
iterations need not monotonically improve the original
unclipped likelihood. This is relevant to the present IRLS draft.

\section{Jensen, the ELBO, and EM monotonicity}
For latent variable $Z$ and any density $q(z)$ supported
where the joint model is positive, Jensen gives
\[
\log p_\theta(y)
=\log\int q(z)\frac{p_\theta(y,z)}{q(z)}\,dz
\ge \mathcal L(q,\theta)
:=\E_q[\log p_\theta(y,Z)]-\E_q[\log q(Z)].
\]
The identity
$\log p_\theta(y)=\mathcal L(q,\theta)
 +D_{\mathrm{KL}}(q\|p_\theta(z\mid y))$
shows exactly when the lower bound is tight.
Exact EM alternates an E-step setting $q$ to the posterior and
an M-step not decreasing the ELBO. This gives nondecreasing
observed-data likelihood, but \emph{not} a global optimum.
For phase--amplitude coupling, the Tort modulation index is
a normalized KL divergence between a binned phase-dependent
amplitude distribution and the uniform distribution; positivity
comes from Gibbs' inequality, while binning and surrogates
are additional statistical decisions.

\section{KKT and alternating least squares}
For differentiable convex $f$ with inequality constraints
$g_j(\beta)\le0$ and equality constraints $h_k(\beta)=0$,
KKT conditions (with appropriate constraint qualifications)
require stationarity, feasibility, nonnegative multipliers and
complementary slackness. Sparse Lasso penalties require
subgradients rather than ordinary derivatives at zero.
For a CP tensor decomposition, exact minimization over a factor
block cannot increase the unpenalized objective \emph{of that
specified block problem}; alternating such steps does not
make the multi-factor objective convex or guarantee global
optimality. The exploratory \path{tensor_cp_als.m}
includes ridge terms and does not track a monotonic objective;
a monotonicity claim for that code must be tested, not assumed.